% Shared exercise chunk: l4-stability
\begin{exercisechunk}
\item \textbf{Contraction and repeated errors.}
\label[exercise]{ex:stability}
This exercise compares recovery from an initial error with continual
correction of new errors.
\begin{enumerate}[(a),beginpenalty=10000]
\item Prove \cref{eq:dobrushin}.
\item For the recovery family \(Q_r\), compute the contraction factor
and use \cref{eq:stability-recursion} to recover the terminal and
count bounds in \cref{prop:recovery}. Verify that the recurrence
inequality is an equality here, giving the exact costs.
\item Now let \(Q_{\varepsilon,r}\), with \(X_0=0\), use
\[
 Q_{\varepsilon,r}(X_{t+1}=1\mid X_t=0)=\varepsilon,\qquad
 Q_{\varepsilon,r}(X_{t+1}=0\mid X_t=1)=r,
\]
where \(0<\varepsilon<1\), \(0\le r\le1\), and \(0\le t<T\).
For \(p_t=Q_{\varepsilon,r}(X_t=1)\), prove
\begin{equation}
 p_0=0,\qquad
 p_{t+1}=\varepsilon+(1-\varepsilon-r)p_t,\qquad
 p_t=\frac{\varepsilon}{\varepsilon+r}
       [1-(1-\varepsilon-r)^t].
 \label{eq:recovery}
\end{equation}
Identify the choices of \(r\) giving the two models in
\cref{ex:binary-count}.
\item Compute the contraction factor \(|1-\varepsilon-r|\).
For \(0\le r\le1-\varepsilon\), prove
\begin{equation}
 \E_{Q_{\varepsilon,r}}[c_{\mathrm{fin}}]\le
 \frac{\varepsilon}{\varepsilon+r},\qquad
 \frac1T\E_{Q_{\varepsilon,r}}[c_{\mathrm{cnt}}]\le
 \frac{\varepsilon}{\varepsilon+r}.
 \label{eq:recovery-cost}
\end{equation}
Recover these bounds from \cref{thm:markov-stability}, using
\(e_1=\varepsilon\) and local error \(\varepsilon\) at every
subsequent step. At \(r=0\), why does the positive contraction margin
still give the bound one? How large must \(r\) be relative to
\(\varepsilon\) for the limiting error probability to be at most
\(s\in(0,1)\)?
\end{enumerate}
\begin{hints}
For (a), normalize the positive and negative parts of
\(\alpha-\beta\) to obtain two probability laws.
\end{hints}
\end{exercisechunk}
